\titlespacing*{\subsection}{0pt}{1.0em}{0.3em}
\titlespacing*{\paragraph}{0pt}{1.0em}{0.3em}

본 연구는 글로벌 외환(FX) 시장에 내재된 국가 간 통화의 비대칭적 상호의존 관계를 시스템 차원에서 규명하기 위해, 그래프 신경망 기반의 예측 프레임워크인 B-MTGNN을 제안하고 그 유효성을 검증하였다. 기존의 단일 시계열 기반 모델 또는 두 통화 간 비교 모델의 한계를 넘어, 다변량 시계열 처리 알고리즘과 그래프 신경망(GNN)을 결합함으로써 통화 간의 복잡한 상호작용을 모델링하였다. 특히, 실증 분석을 통해 개별 국가의 환율 다이내믹은 구조적인 비대칭성을 띠며, 네트워크 상호작용에 강한 의존성을 보임이 확인되었다. 나아가 본 연구는 수치적 예측의 한계를 극복하고자 LangGraph 기반의 멀티 에이전트 시스템을 도입하였다. B-MTGNN의 예측 결과와 불확실성 데이터를 RAG 기술과 결합함으로써 거시경제적 맥락 내에서 설명 가능성을 확보하고, 분석의 완결성을 높였다.
제안된 프레임워크는 수치 예측과 지능형 추론을 통합한 새로운 분석 방법을 제시한다는 점에서 학술적 의의를 갖는다. 연구 결과는 국제 금융시장을 분석할 때 네트워크 구조 및 불확실성 정량화뿐만 아니라, 멀티 에이전트 기반의 맥락적 추론 과정을 병행하는 것이 필수적임을 시사한다. 이러한 다각적 접근법은 향후 비대칭적 상호작용이 존재하는 다양한 금융 도메인으로 확장할 수 있는 모델링 패러다임을 제공할 것으로 기대된다.
